\chapter*{Preface}
Add up the first ten cubes: $1+8+27+\dots+1000=3025$. Now add the first ten numbers and square the answer: $1+2+\dots+10=55$, and $55^2=3025$. The two numbers agree. Try $n=4$ and they agree again: both sides are $100$. Do they agree for every $n$, or only for the ones you happened to try?

Here is a different kind of question. Four guests leave their hats at the door, and the hats come back in random order. How often does nobody get their own hat back? There are $24$ ways to hand back four hats, and in $9$ of them every guest walks off with the wrong one. That is $9/24$, a little over a third. With ten guests, or ten million, the answer barely moves.

\begin{center}
\fig{preface-two-questions}\\[0.35em]
{\small\itshape Left: the first four values agree, and so does $n=10$. Right: one way four guests can each receive someone else's hat.}
\end{center}

The two questions want different kinds of answer. The cubes make a claim about every $n$, and no pile of examples settles that. You can check a million cases and the next one is still unchecked. What settles it is an argument that shows why the agreement cannot stop. The hats make a claim about a proportion. Here nobody asks whether it always happens. The question is how often it happens, and why that fraction settles where it does. The series is called \emph{Order and Chance: The Hidden Architecture of Mathematics}, and those are its two kinds of sentence: patterns that have to hold in every case, and proportions that hold across all the cases together. Each kind needs a proof, and by the end of this book you will have proved both of these.

\section*{What this book is for}
This book teaches you to do three things. You will learn to read a mathematical sentence and say exactly what it claims. You will learn to look for a proof, which usually means trying small cases, making a guess, getting stuck, and finding a better way to say the problem. And you will learn to write the proof down so that someone else can check every line.

You need secondary-school algebra and nothing more: fractions, powers, equations, inequalities, and a little feeling for graphs and elementary geometry. You do not need a previous course in proof, calculus, linear algebra or probability. Every tool is built here before it is used. Some problems will still take you an afternoon. That is the work itself. It does not mean you missed something earlier.

Along the way you will meet induction and binomial coefficients, counting, limits and series, complex numbers and roots of unity, polynomials, linear algebra, calculus, inequalities and convexity, and first probability. Each subject gets one job: to give you a tool you can prove things with. None of them is treated in the depth a specialist course would give it. Later volumes take that up. What you get here is enough of each to use it honestly, and the habit of asking what each step is allowed to assume.

The book is not a survey, and it is not a collection of tricks. A result appears here because you can prove it from what came before, or because the text says plainly that it cannot be proved yet.

\section*{What we take for granted}
Every proof starts somewhere. Euclid's \emph{Elements} opens by listing what its proofs may use. One of the common notions, in Heath's translation, is that things equal to the same thing are equal to one another. Later propositions may use the list. They may not add to it quietly. Here is our list.

We use classical logic, ordinary set notation, and the usual arithmetic and order of the natural numbers $\N=\{0,1,2,\ldots\}$, the integers $\Z$, the rationals $\Q$ and the reals $\R$, with their usual inclusions. The positive integers are $\N^*=\{1,2,\ldots\}$. We assume that $\N$ is well ordered: every nonempty subset has a least element. Chapter~\ref{ch:sums} states that fact and derives induction from it, together with the elementary divisibility facts it needs.

For the reals we assume one more thing, completeness: every nonempty set of reals that is bounded above has a least upper bound. What the book needs about limits comes from that sentence and ordinary algebra. The book proves the convergence tests it uses, the intermediate-value and extreme-value theorems, uniform continuity on a closed bounded interval, Rolle's theorem and the mean value theorem. None of them is assumed separately. A short toolkit of derivatives is in Chapter~\ref{ch:sequences}. The full rules, with proofs, are in Chapter~\ref{ch:calculus}.

We take the functions $\exp$, $\ln$, $\cos$ and $\sin$ as given on their usual real domains, with their addition laws and their derivatives: $\exp'=\exp$, $\ln'(x)=1/x$ for $x>0$, $\sin'=\cos$ and $\cos'=-\sin$. We also assume $\exp(0)=1$, that $\exp$ is a bijection from $\R$ onto the positive reals with inverse $\ln$, and that sine and cosine have the usual unit-circle normalisation and period $2\pi$. When the text uses a power series to compute one of these functions, it proves that the series gives the right value. The series does not stand in for a construction we skipped.

We assume the integral of a continuous function on a closed bounded interval, with four properties: it is linear, it preserves order, it adds over adjacent intervals, and a constant $c$ on $[a,b]$ integrates to $c(b-a)$. From those, the book proves the fundamental theorem of calculus, integration by parts, change of variables and the error estimates it uses. An improper integral is a limit. It does not give permission to swap two infinite processes.

Later volumes build these objects from earlier ground. Nothing from that later work is used here as an unproved lemma.

Probability starts with finite sets of outcomes, where every probability is a finite sum. A few examples with countably many outcomes are marked as an extension, and they use convergent sums and truncation. Building probability on an infinite sequence of coin tosses needs measure theory, and that belongs to a later volume. When a challenge box points past this book, it says which part is proved here, which part is a view ahead, and where to read more.

History comes with sources. When the sums chapter tells the schoolroom story about the young Gauss, it keeps the earliest account apart from what was added later.

\section*{How to read it}
Read with a pencil. When a Pause and think box asks a question, write down your guess before you read on. A wrong guess you made yourself teaches more than a right answer you only read. A workshop shows how a proof is found, not just how it looks at the end: a few examples, an attempt that fails, and the change of viewpoint that shows why it failed.

Exercises come in three levels. One star builds technique and care. Two stars ask you to choose a method or join several steps. Three stars ask for a longer piece of discovery. Every numbered exercise has a hint and a full solution, kept separate. In each solution, \emph{How to find it} shows the search and \emph{Write-up} gives the argument that came out of it. Try the problem before you read either.

\subsection*{Two routes through the book}

The \emph{foundations route} is for a first reading. Read the chapters in order, including the worked examples and the proof workshops. Try the one-star exercises before looking at their hints. In the counting chapter, Exercise~\ref{ex:count-paths-first} is a short stepping stone to the Catalan paths. If a two- or three-star problem stops you for an hour, leave it marked and carry on. The next chapter does not require you to finish every challenge.

The \emph{exploration route} is for a second reading, or for a reader who already knows some of the subject. Return to the two-star exercises and then choose three-star problems that interest you. The stars describe the work required, not the value of a problem. Some of the later bridge problems borrow a theorem from another chapter or volume. Their statements say where that tool comes from. Read the tool first, or treat the problem as a preview; it is not an unannounced prerequisite for the next chapter.

Here is a way to tackle either route. For a count, list the smallest objects and decide whether the cases overlap. For an induction, write the first case and the exact statement you will carry to the next. For a limit, name the quantity you need to make small. For a difficult exercise, try one of these smaller jobs before reading its hint. The solution then shows not just which calculation worked, but how it was found.

The field guide at the end brings these habits together for a problem that does not come with a chapter name.

This is Volume~1, \emph{First Principles}: \emph{Proof, Problem Solving and Mathematical Thinking}. You can read it without any other volume. It starts before the cubes, with a smaller question: what does a mathematical sentence say? The equation $x^2=1$ asks three different questions depending on its domain and the words around it, and the next chapter reads all three. The chapter after that proves the cube identity for every $n$, four different ways. The hats wait for the counting chapter, which counts them exactly, and the probability chapter, which shows why the proportion settles near $1/e$. After this book comes Volume~2, \emph{The Integers: Primes, Patterns and Proofs}. A prime is easy to define and hard to control, and this volume stops before those proofs begin.
